%%%%%%%%%%%%%%%%%%%%%%%%%%%%%%%%%%%%%%%%%
% The Legrand Orange Book
% LaTeX Template
% Version 2.0 (9/2/15)
%
% This template has been downloaded from:
% http://www.LaTeXTemplates.com
%
% Mathias Legrand (legrand.mathias@gmail.com) with modifications by:
% Vel (vel@latextemplates.com)
%
% License:
% CC BY-NC-SA 3.0 (http://creativecommons.org/licenses/by-nc-sa/3.0/)
%
% Compiling this template:
% This template uses biber for its bibliography and makeindex for its index.
% When you first open the template, compile it from the command line with the 
% commands below to make sure your LaTeX distribution is configured correctly:
%
% 1) pdflatex main
% 2) makeindex main.idx -s StyleInd.ist
% 3) biber main
% 4) pdflatex main x 2
%
% After this, when you wish to update the bibliography/index use the appropriate
% command above and make sure to compile with pdflatex several times 
% afterwards to propagate your changes to the document.
%
% This template also uses a number of packages which may need to be
% updated to the newest versions for the template to compile. It is strongly
% recommended you update your LaTeX distribution if you have any
% compilation errors.
%
% Important note:
% Chapter heading images should have a 2:1 width:height ratio,
% e.g. 920px width and 460px height.
%
%%%%%%%%%%%%%%%%%%%%%%%%%%%%%%%%%%%%%%%%%

%----------------------------------------------------------------------------------------
%	PACKAGES AND OTHER DOCUMENT CONFIGURATIONS
%----------------------------------------------------------------------------------------

\documentclass[11pt,fleqn]{book} % Default font size and left-justified equations

%----------------------------------------------------------------------------------------

%%%%%%%%%%%%%%%%%%%%%%%%%%%%%%%%%%%%%%%%%
% The Legrand Orange Book
% Structural Definitions File
% Version 2.0 (9/2/15)
%
% Original author:
% Mathias Legrand (legrand.mathias@gmail.com) with modifications by:
% Vel (vel@latextemplates.com)
% 
% This file has been downloaded from:
% http://www.LaTeXTemplates.com
%
% License:
% CC BY-NC-SA 3.0 (http://creativecommons.org/licenses/by-nc-sa/3.0/)
%
%%%%%%%%%%%%%%%%%%%%%%%%%%%%%%%%%%%%%%%%%

%----------------------------------------------------------------------------------------
%	VARIOUS REQUIRED PACKAGES AND CONFIGURATIONS
%----------------------------------------------------------------------------------------

\usepackage[top=3cm,bottom=3cm,left=3cm,right=3cm,headsep=10pt,a4paper]{geometry} % Page margins

\usepackage{graphicx} % Required for including pictures
\graphicspath{{Pictures/}} % Specifies the directory where pictures are stored

\usepackage{lipsum} % Inserts dummy text

\usepackage{tikz} % Required for drawing custom shapes

\usepackage[english]{babel} % English language/hyphenation

\usepackage{enumitem} % Customize lists
\setlist{nolistsep} % Reduce spacing between bullet points and numbered lists

\usepackage{booktabs} % Required for nicer horizontal rules in tables

\usepackage{xcolor} % Required for specifying colors by name
\definecolor{ocre}{RGB}{243,102,25} % Define the orange color used for highlighting throughout the book

%----------------------------------------------------------------------------------------
%	FONTS
%----------------------------------------------------------------------------------------

\usepackage{avant} % Use the Avantgarde font for headings
%\usepackage{times} % Use the Times font for headings
\usepackage{mathptmx} % Use the Adobe Times Roman as the default text font together with math symbols from the Sym­bol, Chancery and Com­puter Modern fonts

\usepackage{microtype} % Slightly tweak font spacing for aesthetics
\usepackage[utf8]{inputenc} % Required for including letters with accents
\usepackage[T1]{fontenc} % Use 8-bit encoding that has 256 glyphs

%----------------------------------------------------------------------------------------
%	BIBLIOGRAPHY AND INDEX
%----------------------------------------------------------------------------------------

\usepackage[style=alphabetic,citestyle=numeric,sorting=nyt,sortcites=true,autopunct=true,babel=hyphen,hyperref=true,abbreviate=false,backref=true,backend=biber]{biblatex}
\addbibresource{bibliography.bib} % BibTeX bibliography file
\defbibheading{bibempty}{}

\usepackage{calc} % For simpler calculation - used for spacing the index letter headings correctly
\usepackage{makeidx} % Required to make an index
\makeindex % Tells LaTeX to create the files required for indexing

%----------------------------------------------------------------------------------------
%	MAIN TABLE OF CONTENTS
%----------------------------------------------------------------------------------------

\usepackage{titletoc} % Required for manipulating the table of contents

\contentsmargin{0cm} % Removes the default margin

% Part text styling
\titlecontents{part}[0cm]
{\addvspace{20pt}\centering\large\bfseries}
{}
{}
{}

% Chapter text styling
\titlecontents{chapter}[1.25cm] % Indentation
{\addvspace{12pt}\large\sffamily\bfseries} % Spacing and font options for chapters
{\color{ocre!60}\contentslabel[\Large\thecontentslabel]{1.25cm}\color{ocre}} % Chapter number
{\color{ocre}}  
{\color{ocre!60}\normalsize\;\titlerule*[.5pc]{.}\;\thecontentspage} % Page number

% Section text styling
\titlecontents{section}[1.25cm] % Indentation
{\addvspace{3pt}\sffamily\bfseries} % Spacing and font options for sections
{\contentslabel[\thecontentslabel]{1.25cm}} % Section number
{}
{\hfill\color{black}\thecontentspage} % Page number
[]

% Subsection text styling
\titlecontents{subsection}[1.25cm] % Indentation
{\addvspace{1pt}\sffamily\small} % Spacing and font options for subsections
{\contentslabel[\thecontentslabel]{1.25cm}} % Subsection number
{}
{\ \titlerule*[.5pc]{.}\;\thecontentspage} % Page number
[]

% List of figures
\titlecontents{figure}[0em]
{\addvspace{-5pt}\sffamily}
{\thecontentslabel\hspace*{1em}}
{}
{\ \titlerule*[.5pc]{.}\;\thecontentspage}
[]

% List of tables
\titlecontents{table}[0em]
{\addvspace{-5pt}\sffamily}
{\thecontentslabel\hspace*{1em}}
{}
{\ \titlerule*[.5pc]{.}\;\thecontentspage}
[]

%----------------------------------------------------------------------------------------
%	MINI TABLE OF CONTENTS IN PART HEADS
%----------------------------------------------------------------------------------------

% Chapter text styling
\titlecontents{lchapter}[0em] % Indenting
{\addvspace{15pt}\large\sffamily\bfseries} % Spacing and font options for chapters
{\color{ocre}\contentslabel[\Large\thecontentslabel]{1.25cm}\color{ocre}} % Chapter number
{}  
{\color{ocre}\normalsize\sffamily\bfseries\;\titlerule*[.5pc]{.}\;\thecontentspage} % Page number

% Section text styling
\titlecontents{lsection}[0em] % Indenting
{\sffamily\small} % Spacing and font options for sections
{\contentslabel[\thecontentslabel]{1.25cm}} % Section number
{}
{}

% Subsection text styling
\titlecontents{lsubsection}[.5em] % Indentation
{\normalfont\footnotesize\sffamily} % Font settings
{}
{}
{}

%----------------------------------------------------------------------------------------
%	PAGE HEADERS
%----------------------------------------------------------------------------------------

\usepackage{fancyhdr} % Required for header and footer configuration

\pagestyle{fancy}
\renewcommand{\chaptermark}[1]{\markboth{\sffamily\normalsize\bfseries\chaptername\ \thechapter.\ #1}{}} % Chapter text font settings
\renewcommand{\sectionmark}[1]{\markright{\sffamily\normalsize\thesection\hspace{5pt}#1}{}} % Section text font settings
\fancyhf{} \fancyhead[LE,RO]{\sffamily\normalsize\thepage} % Font setting for the page number in the header
\fancyhead[LO]{\rightmark} % Print the nearest section name on the left side of odd pages
\fancyhead[RE]{\leftmark} % Print the current chapter name on the right side of even pages
\renewcommand{\headrulewidth}{0.5pt} % Width of the rule under the header
\addtolength{\headheight}{2.5pt} % Increase the spacing around the header slightly
\renewcommand{\footrulewidth}{0pt} % Removes the rule in the footer
\fancypagestyle{plain}{\fancyhead{}\renewcommand{\headrulewidth}{0pt}} % Style for when a plain pagestyle is specified

% Removes the header from odd empty pages at the end of chapters
\makeatletter
\renewcommand{\cleardoublepage}{
\clearpage\ifodd\c@page\else
\hbox{}
\vspace*{\fill}
\thispagestyle{empty}
\newpage
\fi}

%----------------------------------------------------------------------------------------
%	THEOREM STYLES
%----------------------------------------------------------------------------------------

\usepackage{amsmath,amsfonts,amssymb,amsthm} % For math equations, theorems, symbols, etc

\newcommand{\intoo}[2]{\mathopen{]}#1\,;#2\mathclose{[}}
\newcommand{\ud}{\mathop{\mathrm{{}d}}\mathopen{}}
\newcommand{\intff}[2]{\mathopen{[}#1\,;#2\mathclose{]}}
\newtheorem{notation}{Notation}[chapter]

% Boxed/framed environments
\newtheoremstyle{ocrenumbox}% % Theorem style name
{0pt}% Space above
{0pt}% Space below
{\normalfont}% % Body font
{}% Indent amount
{\small\bf\sffamily\color{ocre}}% % Theorem head font
{\;}% Punctuation after theorem head
{0.25em}% Space after theorem head
{\small\sffamily\color{ocre}\thmname{#1}\nobreakspace\thmnumber{\@ifnotempty{#1}{}\@upn{#2}}% Theorem text (e.g. Theorem 2.1)
\thmnote{\nobreakspace\the\thm@notefont\sffamily\bfseries\color{black}---\nobreakspace#3.}} % Optional theorem note
\renewcommand{\qedsymbol}{$\blacksquare$}% Optional qed square

\newtheoremstyle{blacknumex}% Theorem style name
{5pt}% Space above
{5pt}% Space below
{\normalfont}% Body font
{} % Indent amount
{\small\bf\sffamily}% Theorem head font
{\;}% Punctuation after theorem head
{0.25em}% Space after theorem head
{\small\sffamily{\tiny\ensuremath{\blacksquare}}\nobreakspace\thmname{#1}\nobreakspace\thmnumber{\@ifnotempty{#1}{}\@upn{#2}}% Theorem text (e.g. Theorem 2.1)
\thmnote{\nobreakspace\the\thm@notefont\sffamily\bfseries---\nobreakspace#3.}}% Optional theorem note

\newtheoremstyle{blacknumbox} % Theorem style name
{0pt}% Space above
{0pt}% Space below
{\normalfont}% Body font
{}% Indent amount
{\small\bf\sffamily}% Theorem head font
{\;}% Punctuation after theorem head
{0.25em}% Space after theorem head
{\small\sffamily\thmname{#1}\nobreakspace\thmnumber{\@ifnotempty{#1}{}\@upn{#2}}% Theorem text (e.g. Theorem 2.1)
\thmnote{\nobreakspace\the\thm@notefont\sffamily\bfseries---\nobreakspace#3.}}% Optional theorem note

% Non-boxed/non-framed environments
\newtheoremstyle{ocrenum}% % Theorem style name
{5pt}% Space above
{5pt}% Space below
{\normalfont}% % Body font
{}% Indent amount
{\small\bf\sffamily\color{ocre}}% % Theorem head font
{\;}% Punctuation after theorem head
{0.25em}% Space after theorem head
{\small\sffamily\color{ocre}\thmname{#1}\nobreakspace\thmnumber{\@ifnotempty{#1}{}\@upn{#2}}% Theorem text (e.g. Theorem 2.1)
\thmnote{\nobreakspace\the\thm@notefont\sffamily\bfseries\color{black}---\nobreakspace#3.}} % Optional theorem note
\renewcommand{\qedsymbol}{$\blacksquare$}% Optional qed square
\makeatother

% Defines the theorem text style for each type of theorem to one of the three styles above
\newcounter{dummy} 
\numberwithin{dummy}{section}
\theoremstyle{ocrenumbox}
\newtheorem{theoremeT}[dummy]{Theorem}
\newtheorem{problem}{Problem}[chapter]
\newtheorem{exerciseT}{Exercise}[chapter]
\theoremstyle{blacknumex}
\newtheorem{exampleT}{Example}[chapter]
\theoremstyle{blacknumbox}
\newtheorem{vocabulary}{Vocabulary}[chapter]
\newtheorem{definitionT}{Definition}[section]
\newtheorem{corollaryT}[dummy]{Corollary}
\theoremstyle{ocrenum}
\newtheorem{proposition}[dummy]{Proposition}

%----------------------------------------------------------------------------------------
%	DEFINITION OF COLORED BOXES
%----------------------------------------------------------------------------------------

\RequirePackage[framemethod=default]{mdframed} % Required for creating the theorem, definition, exercise and corollary boxes

% Theorem box
\newmdenv[skipabove=7pt,
skipbelow=7pt,
backgroundcolor=black!5,
linecolor=ocre,
innerleftmargin=5pt,
innerrightmargin=5pt,
innertopmargin=5pt,
leftmargin=0cm,
rightmargin=0cm,
innerbottommargin=5pt]{tBox}

% Exercise box	  
\newmdenv[skipabove=7pt,
skipbelow=7pt,
rightline=false,
leftline=true,
topline=false,
bottomline=false,
backgroundcolor=ocre!10,
linecolor=ocre,
innerleftmargin=5pt,
innerrightmargin=5pt,
innertopmargin=5pt,
innerbottommargin=5pt,
leftmargin=0cm,
rightmargin=0cm,
linewidth=4pt]{eBox}	

% Definition box
\newmdenv[skipabove=7pt,
skipbelow=7pt,
rightline=false,
leftline=true,
topline=false,
bottomline=false,
linecolor=ocre,
innerleftmargin=5pt,
innerrightmargin=5pt,
innertopmargin=0pt,
leftmargin=0cm,
rightmargin=0cm,
linewidth=4pt,
innerbottommargin=0pt]{dBox}	

% Corollary box
\newmdenv[skipabove=7pt,
skipbelow=7pt,
rightline=false,
leftline=true,
topline=false,
bottomline=false,
linecolor=gray,
backgroundcolor=black!5,
innerleftmargin=5pt,
innerrightmargin=5pt,
innertopmargin=5pt,
leftmargin=0cm,
rightmargin=0cm,
linewidth=4pt,
innerbottommargin=5pt]{cBox}

% Creates an environment for each type of theorem and assigns it a theorem text style from the "Theorem Styles" section above and a colored box from above
\newenvironment{theorem}{\begin{tBox}\begin{theoremeT}}{\end{theoremeT}\end{tBox}}
\newenvironment{exercise}{\begin{eBox}\begin{exerciseT}}{\hfill{\color{ocre}\tiny\ensuremath{\blacksquare}}\end{exerciseT}\end{eBox}}				  
\newenvironment{definition}{\begin{dBox}\begin{definitionT}}{\end{definitionT}\end{dBox}}	
\newenvironment{example}{\begin{exampleT}}{\hfill{\tiny\ensuremath{\blacksquare}}\end{exampleT}}		
\newenvironment{corollary}{\begin{cBox}\begin{corollaryT}}{\end{corollaryT}\end{cBox}}	

%----------------------------------------------------------------------------------------
%	REMARK ENVIRONMENT
%----------------------------------------------------------------------------------------

\newenvironment{remark}{\par\vspace{10pt}\small % Vertical white space above the remark and smaller font size
\begin{list}{}{
\leftmargin=35pt % Indentation on the left
\rightmargin=25pt}\item\ignorespaces % Indentation on the right
\makebox[-2.5pt]{\begin{tikzpicture}[overlay]
\node[draw=ocre!60,line width=1pt,circle,fill=ocre!25,font=\sffamily\bfseries,inner sep=2pt,outer sep=0pt] at (-15pt,0pt){\textcolor{ocre}{R}};\end{tikzpicture}} % Orange R in a circle
\advance\baselineskip -1pt}{\end{list}\vskip5pt} % Tighter line spacing and white space after remark

%----------------------------------------------------------------------------------------
%	SECTION NUMBERING IN THE MARGIN
%----------------------------------------------------------------------------------------

\makeatletter
\renewcommand{\@seccntformat}[1]{\llap{\textcolor{ocre}{\csname the#1\endcsname}\hspace{1em}}}                    
\renewcommand{\section}{\@startsection{section}{1}{\z@}
{-4ex \@plus -1ex \@minus -.4ex}
{1ex \@plus.2ex }
{\normalfont\large\sffamily\bfseries}}
\renewcommand{\subsection}{\@startsection {subsection}{2}{\z@}
{-3ex \@plus -0.1ex \@minus -.4ex}
{0.5ex \@plus.2ex }
{\normalfont\sffamily\bfseries}}
\renewcommand{\subsubsection}{\@startsection {subsubsection}{3}{\z@}
{-2ex \@plus -0.1ex \@minus -.2ex}
{.2ex \@plus.2ex }
{\normalfont\small\sffamily\bfseries}}                        
\renewcommand\paragraph{\@startsection{paragraph}{4}{\z@}
{-2ex \@plus-.2ex \@minus .2ex}
{.1ex}
{\normalfont\small\sffamily\bfseries}}

%----------------------------------------------------------------------------------------
%	PART HEADINGS
%----------------------------------------------------------------------------------------

% numbered part in the table of contents
\newcommand{\@mypartnumtocformat}[2]{%
\setlength\fboxsep{0pt}%
\noindent\colorbox{ocre!20}{\strut\parbox[c][.7cm]{\ecart}{\color{ocre!70}\Large\sffamily\bfseries\centering#1}}\hskip\esp\colorbox{ocre!40}{\strut\parbox[c][.7cm]{\linewidth-\ecart-\esp}{\Large\sffamily\centering#2}}}%
%%%%%%%%%%%%%%%%%%%%%%%%%%%%%%%%%%
% unnumbered part in the table of contents
\newcommand{\@myparttocformat}[1]{%
\setlength\fboxsep{0pt}%
\noindent\colorbox{ocre!40}{\strut\parbox[c][.7cm]{\linewidth}{\Large\sffamily\centering#1}}}%
%%%%%%%%%%%%%%%%%%%%%%%%%%%%%%%%%%
\newlength\esp
\setlength\esp{4pt}
\newlength\ecart
\setlength\ecart{1.2cm-\esp}
\newcommand{\thepartimage}{}%
\newcommand{\partimage}[1]{\renewcommand{\thepartimage}{#1}}%
\def\@part[#1]#2{%
\ifnum \c@secnumdepth >-2\relax%
\refstepcounter{part}%
\addcontentsline{toc}{part}{\texorpdfstring{\protect\@mypartnumtocformat{\thepart}{#1}}{\partname~\thepart\ ---\ #1}}
\else%
\addcontentsline{toc}{part}{\texorpdfstring{\protect\@myparttocformat{#1}}{#1}}%
\fi%
\startcontents%
\markboth{}{}%
{\thispagestyle{empty}%
\begin{tikzpicture}[remember picture,overlay]%
\node at (current page.north west){\begin{tikzpicture}[remember picture,overlay]%	
\fill[ocre!20](0cm,0cm) rectangle (\paperwidth,-\paperheight);
\node[anchor=north] at (4cm,-3.25cm){\color{ocre!40}\fontsize{220}{100}\sffamily\bfseries\@Roman\c@part}; 
\node[anchor=south east] at (\paperwidth-1cm,-\paperheight+1cm){\parbox[t][][t]{8.5cm}{
\printcontents{l}{0}{\setcounter{tocdepth}{1}}%
}};
\node[anchor=north east] at (\paperwidth-1.5cm,-3.25cm){\parbox[t][][t]{15cm}{\strut\raggedleft\color{white}\fontsize{30}{30}\sffamily\bfseries#2}};
\end{tikzpicture}};
\end{tikzpicture}}%
\@endpart}
\def\@spart#1{%
\startcontents%
\phantomsection
{\thispagestyle{empty}%
\begin{tikzpicture}[remember picture,overlay]%
\node at (current page.north west){\begin{tikzpicture}[remember picture,overlay]%	
\fill[ocre!20](0cm,0cm) rectangle (\paperwidth,-\paperheight);
\node[anchor=north east] at (\paperwidth-1.5cm,-3.25cm){\parbox[t][][t]{15cm}{\strut\raggedleft\color{white}\fontsize{30}{30}\sffamily\bfseries#1}};
\end{tikzpicture}};
\end{tikzpicture}}
\addcontentsline{toc}{part}{\texorpdfstring{%
\setlength\fboxsep{0pt}%
\noindent\protect\colorbox{ocre!40}{\strut\protect\parbox[c][.7cm]{\linewidth}{\Large\sffamily\protect\centering #1\quad\mbox{}}}}{#1}}%
\@endpart}
\def\@endpart{\vfil\newpage
\if@twoside
\if@openright
\null
\thispagestyle{empty}%
\newpage
\fi
\fi
\if@tempswa
\twocolumn
\fi}

%----------------------------------------------------------------------------------------
%	CHAPTER HEADINGS
%----------------------------------------------------------------------------------------

\newcommand{\thechapterimage}{}%
\newcommand{\chapterimage}[1]{\renewcommand{\thechapterimage}{#1}}%
\def\@makechapterhead#1{%
{\parindent \z@ \raggedright \normalfont
\ifnum \c@secnumdepth >\m@ne
\if@mainmatter
\begin{tikzpicture}[remember picture,overlay]
\node at (current page.north west)
{\begin{tikzpicture}[remember picture,overlay]
\node[anchor=north west,inner sep=0pt] at (0,0) {\includegraphics[width=\paperwidth]{\thechapterimage}};
\draw[anchor=west] (\Gm@lmargin,-9cm) node [line width=2pt,rounded corners=15pt,draw=ocre,fill=white,fill opacity=0.5,inner sep=15pt]{\strut\makebox[22cm]{}};
\draw[anchor=west] (\Gm@lmargin+.3cm,-9cm) node {\huge\sffamily\bfseries\color{black}\thechapter. #1\strut};
\end{tikzpicture}};
\end{tikzpicture}
\else
\begin{tikzpicture}[remember picture,overlay]
\node at (current page.north west)
{\begin{tikzpicture}[remember picture,overlay]
\node[anchor=north west,inner sep=0pt] at (0,0) {\includegraphics[width=\paperwidth]{\thechapterimage}};
\draw[anchor=west] (\Gm@lmargin,-9cm) node [line width=2pt,rounded corners=15pt,draw=ocre,fill=white,fill opacity=0.5,inner sep=15pt]{\strut\makebox[22cm]{}};
\draw[anchor=west] (\Gm@lmargin+.3cm,-9cm) node {\huge\sffamily\bfseries\color{black}#1\strut};
\end{tikzpicture}};
\end{tikzpicture}
\fi\fi\par\vspace*{270\p@}}}

%-------------------------------------------

\def\@makeschapterhead#1{%
\begin{tikzpicture}[remember picture,overlay]
\node at (current page.north west)
{\begin{tikzpicture}[remember picture,overlay]
\node[anchor=north west,inner sep=0pt] at (0,0) {\includegraphics[width=\paperwidth]{\thechapterimage}};
\draw[anchor=west] (\Gm@lmargin,-9cm) node [line width=2pt,rounded corners=15pt,draw=ocre,fill=white,fill opacity=0.5,inner sep=15pt]{\strut\makebox[22cm]{}};
\draw[anchor=west] (\Gm@lmargin+.3cm,-9cm) node {\huge\sffamily\bfseries\color{black}#1\strut};
\end{tikzpicture}};
\end{tikzpicture}
\par\vspace*{270\p@}}
\makeatother

%----------------------------------------------------------------------------------------
%	HYPERLINKS IN THE DOCUMENTS
%----------------------------------------------------------------------------------------

\usepackage{hyperref}
\hypersetup{hidelinks,backref=true,pagebackref=true,hyperindex=true,colorlinks=false,breaklinks=true,urlcolor= ocre,bookmarks=true,bookmarksopen=false,pdftitle={Title},pdfauthor={Author}}
\usepackage{bookmark}
\bookmarksetup{
open,
numbered,
addtohook={%
\ifnum\bookmarkget{level}=0 % chapter
\bookmarksetup{bold}%
\fi
\ifnum\bookmarkget{level}=-1 % part
\bookmarksetup{color=ocre,bold}%
\fi
}
} % Insert the commands.tex file which contains the majority of the structure behind the template



%%agregué


\usepackage[hang, small,labelfont=bf,up,textfont=it,up]{caption} % Custom captions under/above floats in tables or figures
\usepackage{booktabs} % Horizontal rules in tables
\usepackage{float} % Required for tables and figures in the multi-column environment - they




\usepackage{graphicx} % paquete que permite introducir imágenes

\usepackage{booktabs} % Horizontal rules in tables
\usepackage{float} % Required for tables and figures in the multi-column environment - they

\numberwithin{equation}{section} % Number equations within sections (i.e. 1.1, 1.2, 2.1, 2.2 instead of 1, 2, 3, 4)
\numberwithin{figure}{section} % Number figures within sections (i.e. 1.1, 1.2, 2.1, 2.2 instead of 1, 2, 3, 4)
\numberwithin{table}{section} % Number tables within sections (i.e. 1.1, 1.2, 2.1, 2.2 instead of 1, 2, 3, 4)


\setlength\parindent{0pt} % Removes all indentation from paragraphs - comment this line for an assignment with lots of text

%%hasta aquí


\begin{document}

%----------------------------------------------------------------------------------------
%	TITLE PAGE
%----------------------------------------------------------------------------------------

\begingroup
\thispagestyle{empty}
\begin{tikzpicture}[remember picture,overlay]
\coordinate [below=12cm] (midpoint) at (current page.north);
\node at (current page.north west)
{\begin{tikzpicture}[remember picture,overlay]
\node[anchor=north west,inner sep=0pt] at (0,0) {\includegraphics[width=\paperwidth]{background}}; % Background image
\draw[anchor=north] (midpoint) node [fill=ocre!30!white,fill opacity=0.6,text opacity=1,inner sep=1cm]{\Huge\centering\bfseries\sffamily\parbox[c][][t]{\paperwidth}{\centering Modelos con Variable Dependiente Discreta\\[15pt] % Book title
{\Large Econometría II}\\[20pt] % Subtitle
{\huge Alarcón Castillo Henry\\Champa Del Valle Katherine\\Mayhuasca Gutierrez Victor\\ Bautista Ramos Luis}}}; % Author name
\end{tikzpicture}};
\end{tikzpicture}
\vfill
\endgroup


%----------------------------------------------------------------------------------------
%	COPYRIGHT PAGE
%----------------------------------------------------------------------------------------

%\newpage
%~\vfill
%\thispagestyle{empty}

%\noindent Copyright \copyright\ 2013 John Smith\\ % Copyright notice

%\noindent \textsc{Published by Publisher}\\ % Publisher

%\noindent \textsc{book-website.com}\\ % URL

%\noindent Licensed under the Creative Commons Attribution-NonCommercial 3.0 Unported License (the ``License''). You may not use this file except in compliance with the License. You may obtain a copy of the License at \url{http://creativecommons.org/licenses/by-nc/3.0}. Unless required by applicable law or agreed to in writing, software distributed under the License is distributed on an \textsc{``as is'' basis, without warranties or conditions of any kind}, either express or implied. See the License for the specific language governing permissions and limitations under the License.\\ % License information

%\noindent \textit{First printing, March 2013} % Printing/edition date

%----------------------------------------------------------------------------------------
%	TABLE OF CONTENTS
%----------------------------------------------------------------------------------------

\chapterimage{ima1} % Table of contents heading image

%\chapterimage{chapter_head_1.pdf} % Table of contents heading image

\pagestyle{empty} % No headers

 \tableofcontents % Print the table of contents itself

\cleardoublepage % Forces the first chapter to start on an odd page so it's on the right

\pagestyle{fancy} % Print headers again

%----------------------------------------------------------------------------------------
%	PART
%----------------------------------------------------------------------------------------

\part{Parte Uno}

%----------------------------------------------------------------------------------------
%	CHAPTER 1
%----------------------------------------------------------------------------------------

\chapterimage{ima2} % Chapter heading image

\chapter{Installation}

EMDO101 classic is directly assembled on utility meter on D0 Interface (IEC-62056-21). 
One RS485
one relay output driver
one s0 optical input 
two s0 electrical inputs
two s0 electrical outputs
board temperature sensor
Ethernet 
usb for firmware update
USB Power supply
default reset button
picture emdo101

pinout table
1 S0 IN 1 +
2 S0 IN 1 -
3 S0 IN 2 +
4 S0 IN 2 -
5 S0 OUT 1 +
6 S0 OUT 1 -
7 S0 OUT 2 +
8 S0 OUT 2 -
9 GND
10 5V Out 200mA
11 RS485 A
12 RS485 B
13 Relay OUT

EMDO102 is mounted on DIN Rail and can be connected with an optional external D0 reader head IEC-62056-21
two rs485
two relay output drivers
two s0 electrical inputs
two s0 electrical outputs
board temperature sensor
1 wire
Ethernet 
usb for firmware update
default reset button
Requires 24V DIN rail power supply
picture emdo102

pinout table
1 24V In (or 12V)
2 GND
3 GND
4 24V Out (internal fuse protected, approx. 1.4A)
5 Relay OUT 1
6 Relay OUT 2
7 GND (1Wire)
8 1Wire

1 S0 OUT 1 +
2 S0 OUT 1 -
3 S0 OUT 2 +
4 S0 OUT 2 -
5 5V Out (approx 0.9A in total)
6 GND
7 RS485 A1
8 RS485 B1

1 S0 IN 1 +
2 S0 IN 1 -
3 S0 IN 2 +
4 S0 IN 2 -
5 5V Out (approx 0.9A in total)
6 GND
7 RS485 A2
8 RS485 B2



EMDO103 is connected over MBUS bus
one mbus slave port
one rs485
two relay output drivers
two s0 electrical inputs
two s0 electrical outputs
board temperature sensor
1 wire
Ethernet 
usb for firmware update
default reset button
Requires 24V DIN rail power supply
picture emdo103
pinout table
1 24V In (or 12V)
2 GND
3 GND
4 24V Out (internal fuse protected, approx. 1.4A)
5 Relay OUT 1
6 Relay OUT 2
7 GND (1Wire)
8 1Wire

1 S0 OUT 1 +
2 S0 OUT 1 -
3 S0 OUT 2 +
4 S0 OUT 2 -
5 5V Out (approx 0.9A in total)
6 GND
7 RS485 A1
8 RS485 B1

1 S0 IN 1 +
2 S0 IN 1 -
3 S0 IN 2 +
4 S0 IN 2 -
5 5V Out (approx 0.9A in total)
6 GND
7 MBUS A (requires external power supply over bus)
8 MBUS B (requires external power supply over bus)



EMDO104 is connected over KNX bus
one knx slave port
one rs485
two relay output drivers
two s0 electrical inputs
two s0 electrical outputs
board temperature sensor
1 wire
Ethernet 
usb for firmware update
default reset button
Requires 24V DIN rail power supply
picture emdo104
pinout table
1 24V In (or 12V)
2 GND
3 GND
4 24V Out (internal fuse protected, approx. 1.4A)
5 Relay OUT 1
6 Relay OUT 2
7 GND (1Wire)
8 1Wire

1 S0 OUT 1 +
2 S0 OUT 1 -
3 S0 OUT 2 +
4 S0 OUT 2 -
5 5V Out (approx 0.9A in total)
6 GND
7 RS485 A1
8 RS485 B1

1 S0 IN 1 +
2 S0 IN 1 -
3 S0 IN 2 +
4 S0 IN 2 -
5 5V Out (approx 0.9A in total)
6 GND
7 KNX A (requires external power supply over bus)
8 KNX B (requires external power supply over bus)



Expansion boards
Expansion board for LoRaWAN, WMBUS, Enocean, Zigbee, WiFi, picture

Install / deinstall Expansion Board, disconnect power, open EMDO enclosure with two black skews or remove cover
picture emdo classic skews to open
picture din rail enclosure open
put expansion board on header
for radio modules it is important that electromagnetic radiation is not obscured by walls, metal enclosures,... this greatly disgrades performance.


%----------------------------------------------------------------------------------------
%	CHAPTER 2
%----------------------------------------------------------------------------------------

\chapter{Webinterface}

EMDO runs with dhcp server, however if no dhcp server is present, fixed ip can be used.
if settings are bad and emdo can't be connected anymore with fixed ip, there is a reset button to factory default on the emdo

add pictures from web folderdescribe all pictures, except credits, will go into its own chapter
pictures must be cut on bottom and top that can be done in latex too. pls make pictures again before final manual, we will still change one or the other mask slightly. e.g. tab menu has changed (temp removed).you might take some text from welcome tab.

contact tab: there is support page in case of trouble

network tab (top down): enable dhcp client, set network static settings, ntp server ip or dn , mqtt ip or domain name, remote syslog (rsyslog) server ip or domain name, configure logging, configure remote shell (rsh) should be disabled by default, config tcp to d0 gateway, config tcp to rs485 gateway, auth password.
All data loggin is based on ntp time, thus this is very important.

firmware update tab: System log gives details about system. describe procedure for firmware update from file. Prior to update, make sure data from emdo is copied with get data snapshot. then choose firmware file, there are two files one for program memory on the cpu and second for filesystem. This order should be maintained for firmware updates. 

D0, RS485, S0 Input, S0 Output tab: describe the settings. D0 is for the IEC interface. These settings can also be set from basic interpreter. 

Editor tab: this is the main tab for writing and running a program.  Pls play around with the editor and  describe the buttons. Time is from ntp server, it is not running info is shown. white area with numbers for the source code, blue background area is for print output, variable area to show variable values, buttons to run stop, restart, slow step, upload download the basic interpreter code. it is also possible to upload libraries.
% FIXME add tabs for mbus and knx when available

graphs tab: all data are saved on 32mb flash that is build into the emdo. logger library write data to flash and it is visualized as plots. live data are shown as donuts with text. these data can be configured.

-----------------------------

\chapter{Basic Interpreter}

reference to mmbasic
table with commands
pls make sure all of these are convered.most should be covered already in the word document. all otheres I can quickly write up for you, once it is in table format. to describe some operator we can take text from mmbasic manual

\begin{table}[]
\centering
\caption{Predefined Read Only Variables}
\label{Predefined_Read_Only_Variables}
\begin{tabular}{|p{4cm}|p{10cm}|}
\hline
\textbf{Basic Interpreter} & \textbf{Description}                                                             \\ \hline
MM.VER & The version number of the firmware in the form aa.bbcc where aa is the major version number, bb is the minor version number and cc is the revision number (normally zero but A = 01, B = 02, etc).
\\ \hline
MM.FNAME\$ & The name of the file that will be used as the default for the SAVE command. This is set by LOAD, RUN and SAVE.
\\ \hline
MM.CMDLINE\$ & The command line used with the implied RUN command. See the implied RUN command at the start of the next page for the details.
\\ \hline
MM.ERRNO & Is set to the error number if a statement involving the NOR flash fails or
zero if the operation succeeds. This value is not supported yet.
\\ \hline
\end{tabular}
\end{table}

Square brackets indicate that the parameter or characters are optional.
\begin{table}[]
\centering
\caption{Commands}
\label{Commands}
\begin{tabular}{|p{4cm}|p{10cm}|}
\hline
\textbf{Basic Interpreter} & \textbf{Description}                                                             \\ \hline
‘ (single quotation mark) & Starts a comment and any text following it will be ignored. Comments
can be placed anywhere on a line.
\\ \hline
? (question mark) & Shortcut for the PRINT command.
\\ \hline
\\ \hline
CHAIN file\$ &
Clear the current program from memory, load the new program ('file\$')
into memory and run it starting with the first line.
Unlike the RUN command this command retains the current state of the
program (ie, the value of variables, open files, loaded fonts, open COM
ports, etc). The only exception is any open interrupts which will be
automatically closed.
One program can CHAIN to another which can then chain to another (or
back to the original) an unlimited number of times. As long as a
program can be broken down into modules this command allows
programs of almost unlimited size to be run, even with limited memory.
Communication between the modules can be accomplished by assigning
values to one or more variables which then can be examined by the new
chained program.
Note that another way of squeezing a large program into limited memory
is to use the LIBRARY command.
\\ \hline
CLEAR & Delete all variables and recover the memory used by them.
See ERASE for deleting specific array variables.
\\ \hline
CLOSE [\#]nbr [,[\#]nbr] … &
Close the file(s) previously opened with the file number ‘nbr’.
The \# is optional. Also see the OPEN command.
\\ \hline
CONTINUE & Resume running a program that has been stopped by an END statement,
an error, or CTRL-C. The program will restart with the next statement
following the previous stopping point.
\\ \hline
COPY src\$ TO dest\$ & Copy the file named 'src\$' to another file named 'dest\$'.
\\ \hline
DATA constant[,constant]... & Stores numerical and string constants to be accessed by READ.
In general string constants should be surrounded by double quotes (").
An exception is when the string consists of just alphanumeric characters
that do not represent MMBasic keywords (such as THEN, WHILE, etc).
In that case quotes are not needed.
Numerical constants can also be expressions such as 5 * 60.
\\ \hline
DIM var(dim) , [var(dim)]...
or
DIM var\$(dim) LENGTH n
Examples:
DIM nbr(50)
DIM str\$(20)
DIM a(5,5,5), b(1000)
DIM str\$(200) LENGTH 20
&
Specifies a variable that is an array with one or more dimensions. The
variables can be numbers or strings with multiple declarations separated
by commas.
'dim' is a bracketed list of numbers separated by commas. Each number
specifies the number of elements in each dimension. Normally the
numbering of each dimension starts at 0 but the OPTION BASE
command can be used to change this to 1.
For example: DIM nbr(10, 20) specifies a two dimensional array with
11 elements (0 to 10) in the first dimension and 21 (0 to 20) in the
second dimension. The total number of elements is 231 and because
each number requires 4 bytes a total of 924 bytes of memory will be
allocated.
String arrays will by default be allocated memory for 255 characters for
each element and this can quickly use up memory. In that case the
LENGTH keyword can be used to specify the amount of memory to be
allocated to each element. This allocation ('n') can be from 1 to 255
characters.
For example: DIM str\$(5, 10) will declare a string array with 66
elements consuming 16,896 bytes of memory while:
DIM str\$(5, 10) LENGTH 20
Will only consume 1,386 bytes of memory. Note that the amount of
memory allocated for each element is n + 1 as the extra byte is used to
track the actual length of the string stored in each element.
If a string longer than 'n' is assigned to an element of the array an error
will be produced. Other than this string arrays created with the
LENGTH keyword act exactly the same as other string arrays.
\\ \hline
DO
<statements>
LOOP
& This structure will loop forever; the EXIT command can be used to
terminate the loop or control must be explicitly transferred outside of the
loop by commands like GOTO or RETURN (if in a subroutine).
\\ \hline
DO WHILE expression
<statements>
LOOP
DO
<statements>
LOOP UNTIL expression
& 
Loops while "expression" is true (this is equivalent to the older WHILE-
WEND loop, also implemented in MMBasic). If, at the start, the
expression is false the statements in the loop will not be executed, even
once.
Loops until the expression following UNTIL is true. Because the test is
made at the end of the loop the statements inside the loop will be
executed at least once, even if the expression is false.
\\ \hline
ELSE & Introduces a default condition in a multiline IF statement.
See the multiline IF statement for more details.
\\ \hline
ELSEIF expression THEN & Introduces a secondary condition in a multiline IF statement.
See the multiline IF statement for more details.
\\ \hline
ENDIF & Terminates a multiline IF statement.
See the multiline IF statement for more details.
\\ \hline
END & End the running program and return to the command prompt.
\\ \hline
END FUNCTION & Marks the end of a user defined function. See the FUNCTION
command.
Each sub must have one and only one matching END FUNCTION
statement. Use EXIT FUNCTION if you need to return from a
subroutine from within its body.
Only one space is allowed between END and FUNCTION.
\\ \hline
END SUB & Marks the end of a user defined subroutine. See the SUB command.
Each sub must have one and only one matching END SUB statement.
Use EXIT SUB if you need to return from a subroutine from within its
body.
Only one space is allowed between END and SUB.
\\ \hline
ERASE variable
[,variable]...
&
Deletes arrayed variables and frees up the memory.
Use CLEAR to delete all variables including all arrayed variables.
\\ \hline
EXIT
EXIT DO
EXIT FOR
EXIT FUNCTION
EXIT SUB
 & EXIT by itself or EXIT DO provides an early exit from a DO...LOOP
EXIT FOR provides an early exit from a FOR...NEXT loop.
EXIT FUNCTION provides an early exit from a defined function.
EXIT SUB provides an early exit from a defined subroutine.
Only one space is allowed between the two words.
\\ \hline
FILES [fspec\$] & Lists files in the current directory on the NOR flash.
Question marks
(?) will match any character and an asterisk (*) will match any number
of characters. If omitted, all files will be listed. For example:
*.*
Find all entries
*.TXT
Find all entries with an extension of TXT
E*.*
Find all entries starting with E
X?X.*
Find all three letter file names starting and ending with
X
\\ \hline
FOR counter = start TO
finish [STEP increment]
& Initiates a FOR-NEXT loop with the 'counter' initially set to 'start' and
incrementing in 'increment' steps (default is 1) until 'counter' equals
'finish'. The ‘increment’ must be an integer, but may be negative.
See also the NEXT command.
\\ \hline
FUNCTION xxx (arg1
[,arg2, …])
<statements>
<statements>
xxx = <return value>
END FUNCTION
& Defines a callable function. This is the same as adding a new function to
MMBasic while it is running your program.
'xxx' is the function name and it must meet the specifications for naming
a variable. 'arg1', 'arg2', etc are the arguments or parameters to the
function.
To set the return value of the function you assign the value to the
function's name. For example:
FUNCTION SQUARE(a)
SQUARE = a * a
END FUNCTION
Every definition must have one END FUNCTION statement. When this
is reached the function will return its value to the expression from which
it was called. The command EXIT FUNCTION can be used for an early
exit.
You use the function by using its name and arguments in a program just
as you would a normal MMBasic function. For example:
PRINT SQUARE(56.8)
When the function is called each argument in the caller is matched to the
argument in the function definition. These arguments are available only
inside the function.
Functions can be called with a variable number of arguments. Any
omitted arguments in the function's list will be set to zero or a null string.
Arguments in the caller's list that are a variable (ie, not an expression or
constant) will be passed by reference to the function. This means that
any changes to the corresponding argument in the function will also be
copied to the caller's variable.
You must not jump into or out of a function using commands like
GOTO, GOSUB, etc. Doing so will have undefined side effects
including the possibility of ruining your day.
\\ \hline
GOSUB target & Initiates a subroutine call to the target, which can be a line number or a
label. The subroutine must end with RETURN.
\\ \hline
GOTO target & Branches program execution to the target, which can be a line number or
a label.
\\ \hline
IF expr THEN statement
or
IF expr THEN statement
ELSE statement
 & Evaluates the expression ‘expr' and performs the THEN statement if it is
true or skips to the next line if false. The optional ELSE statement is the
reverse of the THEN test. This type of IF statement is all on one line.
The ‘THEN statement’ construct can be also replaced with:
GOTO linenumber | label’.
\\ \hline
IF expression THEN
<statements>
[ELSE
<statements>]
[ELSEIF expression THEN
<statements>]
ENDIF
& The ELSEIF statement (if present) is executed if the previous condition
is false and it starts a new IF chain with further ELSE and/or ELSEIF
statements as required.
One ENDIF is used to terminate the multiline IF.
\\ \hline
INPUT ["prompt string\$";]
list of variables
& Allows input from the keyboard to a list of variables. The input
command will prompt with a question mark (?).
The input must contain commas to separate each data item if there is
more than one variable.
For example, if the command is: INPUT a, b, c
And the following is typed on the keyboard: 23, 87, 66
Then a = 23 and b = 87 and c = 66
If the "prompt string\$" is specified it will be printed before the question
mark. If the prompt string is terminated with a comma (,) rather than the
semicolon (;) the question mark will be suppressed.
\\ \hline
INPUT \#nbr,
list of variables
& Same as above except that the input is read from a file previously opened
for INPUT as ‘nbr’. See the OPEN command.
\\ \hline
KILL file\$ & Deletes the file specified by ‘file\$’. If there is an extension it must be
specified.
If 'file' is a string constant quote marks around it are optional at the
command prompt but must be specified if the command is used within a
program. Example: KILL “SAMPLE.DAT”
\\ \hline
LET variable = expression & Assigns the value of 'expression' to the variable. LET is automatically
assumed if a statement does not start with a command.
\\ \hline
LIBRARY LOAD \$file
and
LIBRARY UNLOAD \$file
& Will load a library file ('file\$') into general memory. Any user defined
commands and subroutines in the file will then become available to the
running program. Up to 8 files can be loaded simultaneously.
A library file is like any other MMBasic program, with the exception
that any programming code outside the user defined commands and
subroutines in the file will be ignored. The library is not visible to the
user (it is not listed by the LIST command) so it should be tested and
debugged as a normal basic program first. Normally a library file has
the extension ".LIB" and that extension will be automatically added if
the file name ('\$file') does not include the extension.
Libraries can be loaded and unloaded in any order. Libraries can be
loaded from within other libraries and nested to an unlimited extent.
Any library file can be unloaded from memory and the memory returned
to the general pool by using the LIBRARY UNLOAD command.
Library files must not be unloaded from within a library that is currently
being used by the program (the results are undefined but it may crash
MMBasic and cause your hair to go prematurely grey).
This command can be used to load specialised libraries to extend the
functionality of MMBasic. Examples include device drivers, libraries
that provide bit manipulation and libraries of specialised mathematical
functions. The library file is only loaded on the first load command
encountered so it is acceptable to put the same load command into every
part of the program or every subroutine that may need the library.
Another use of the LIBRARY command is to extend the amount of
memory available to a program by only loading sections of code as
needed and then unloading them when their task is finished so that
another function can be loaded.
To prevent fragmentation of memory, functions that use a lot of memory
(like arrays) should be declared first before any
libraries are loaded and then unloaded. This is not supported yet.
\\ \hline
LIST
LIST line
LIST -lastline
LIST firstline -
LIST firstline - lastline
& Lists all lines in a program line or a range of lines.
If –lastline is used it will start with the first line in the program. If
startline- is used it will list to the end of the program.
\\ \hline
LOAD file\$
& Loads a program called ‘file\$’ from the current drive into program
memory.
If an extension is not specified “.BAS” will be added to the file name.
\\ \hline
LOCAL variable [,
variables]
& Defines a list of variable names as local to the subroutine or function.
'variable' can be an array and the array will be dimensioned just as if the
DIM command had been used.
A local variable will only be visible within the procedure and will be
deleted (and the memory reclaimed) when the procedure returns. If the
local variable has the same name as a global variable (used before any
subroutines or functions were called) the global variable will be hidden
by the local variable while the procedure is executed.
\\ \hline
LOOP [UNTIL expression] & Terminates a program loop: see DO.
\\ \hline
NAME old\$ AS new\$ & Rename a file or a directory from ‘old\$’ to ‘new\$’
Unlike the other commands that work with file names the NAME
command cannot accept a full pathname (with directories).
\\ \hline
NEW & Deletes the program in memory and clears all variables.
\\ \hline
NEXT [counter-variable] [, counter-variable], etc & NEXT comes at the end of a FOR-NEXT loop; see FOR. The ‘counter-variable’ specifies exactly which loop is being operated on.
If no ‘counter-variable’ is specified the NEXT will default to the
innermost loop. It is also possible to specify multiple counter-variables
as in:
NEXT x, y, z
\\ \hline
ON nbr GOTO | GOSUB
target[,target, target,...]
& ON either branches (GOTO) or calls a subroutine (GOSUB) based on
the rounded value of 'nbr'; if it is 1, the first target is called, if 2, the
second target is called, etc. Target can be a line number or a label.
\\ \hline
OPEN fname\$ FOR mode
AS [\#]fnbr
& Opens a file for reading or writing.
‘fname’ is the filename (8 chars max) with an optional extension (3 chars
max) separated by a dot (.). It can be prefixed with a directory path. For
example: "DIR1/DIR2/FILE.EXT".
‘mode’ is INPUT, OUTPUT, APPEND or RANDOM.
INPUT will open the file for reading and throw an error if the file does
not exist. OUTPUT will open the file for writing and will automatically
overwrite any existing file with the same name.
APPEND will also open the file for writing but it will not overwrite an
existing file; instead any writes will be appended to the end of the file. If
there is no existing file the APPEND mode will act the same as the
OUTPUT mode (i.e. the file is created then opened for writing). Note:
APPEND is not supported on the flash file system (drive A:).
RANDOM will open the file for both read and write and will allow
random access using the SEEK command. When opened the read/write
pointer is positioned at the end of the file. See Appendix I for more
details.
‘fnbr’ is the file number (1 to 10). The \# is optional. Up to 10 files can
be open simultaneously. The INPUT, LINE INPUT, PRINT, WRITE
and CLOSE commands as well as the EOF() and INPUT\$() functions all
use ‘fnbr’ to identify the file being operated on.
See also OPTION ERROR and MM.ERRNO for error handling.
\\ \hline
OPTION BASE 0
or
OPTION BASE 1
& Set the lowest value for array subscripts to either 0 or 1. The default is
0.
This must be used before any arrays are declared.
\\ \hline
OPTION BREAK nn
& Set the value of the break key to 'nn'. This key is used to interrupt a
running program.
The value of the break key is set to CTRL-C key at startup but it can be
changed to any keyboard key using this command (for example,
OPTION BREAK 156 will set the break key to the F12 key).
Setting this option to an invalid key value (for example, 255) will disable
the break function entirely.
\\ \hline
OPTION ERROR
CONTINUE
or
OPTION ERROR ABORT
 & Set the treatment for errors in file input/output. The option CONTINUE
will cause MMBasic to ignore file related errors. The program must
check the variable MM.ERRNO to determine if and what error has
occurred.
The option ABORT sets the normal behaviour (ie, stop the program and
print an error message). The default is ABORT.
Note that this option only relates to errors reading from or writing to
NOR flash. It does not affect the handling of syntax and other
program errors.
\\ \hline
OPTION PROMPT string\$ & Sets the command prompt to the contents of ‘string\$’ (which can also be
an expression which will be evaluated when the prompt is printed).
For example:
OPTION PROMPT “Ok “
or
OPTION PROMPT TIME\$ + “: “
or
OPTION PROMPT CWD\$ + “: “
Maximum length of the prompt string is 48 characters. The prompt is
reset to the default (“> “) on power up but you can automatically set it by
saving the following example program as “AUTORUN.BAS” on the
NOR flash:
10 OPTION PROMPT “My prompt: “
20 NEW
\\ \hline
PAUSE delay & Halt execution of the running program for ‘delay’ mS.
The maximum delay is 2147483647 mS (about 24 days).
\\ \hline
POKE hiword, loword, val
or
POKE keyword, offset, val
or
POKE VAR var, +-offset, val
 & Will set a byte within the PIC32 virtual memory space.
The address is specified by ‘hiword’ which is the top 16 bits of the
address while ‘loword’ is the bottom 16 bits.
Alternatively 'keyword' can be used and 'offset' is the +-offset from the
address of the keyword. The keyword can be
PROGMEM (program memory) or VARTBL (the variable table).
You can also access the memory allocated to a variable by using the
variable's name ('var') preceded by the keyword VAR. This can be used
to access the individual bytes of a numeric variable or a large segment of
RAM allocated to an array. The first element of an array (eg, nbr(0)) is
the start of RAM allocated to the whole array. For example:
DIM nbr(1024)
'allocate 4KB
POKE VAR nbr(0),100,1 'set the 100th byte to 1
This command is for expert users only. address space. Note that MMBasic stores most data (including video) as
32 bit integers and the PIC32 uses little endian format.
WARNING: No validation of the parameters is made and if you use this
facility to access an invalid memory address you will get an "internal
error" which causes the processor to reset and clear all memory.
\\ \hline
PRINT expression
[[,; ]expression] … etc
& Outputs text to the screen. Multiple expressions can be used and must be
separated by either a:
? Comma (,) which will output the tab character
? Semicolon (;) which will not output anything (it is just used to
separate expressions).
? Nothing or a space which will act the same as a semicolon.
A semicolon (;) at the end of the expression list will suppress the
automatic output of a carriage return/ newline at the end of a print
statement.
When printed, a number is preceded with a space if positive or a minus
(-) if negative but is not followed by a space. Integers (whole numbers)
are printed without a decimal point while fractions are printed with the
decimal point and the significant decimal digits. Large numbers (greater
than six digits) are printed in scientific format.
The function FORMAT\$() can be used to format numbers. The function
TAB() can be used to space to a certain column and the string functions
can be used to justify or otherwise format strings.
\\ \hline
QUIT & Will exit MMBasic and return control to the operating system.
\\ \hline
RANDOMIZE nbr & Seed the random number generator with ‘nbr’.
On power up the random number generator is seeded with zero and will
generate the same sequence of random numbers each time. To generate
a different random sequence each time you must use a different value for
‘nbr’. One good way to do this is use the TIMER function.
For example  100 RANDOMIZE TIMER
\\ \hline
READ variable[, variable]...
& Reads values from DATA statements and assigns these values to the
named variables. Variable types in a READ statement must match the
data types in DATA statements as they are read. See also DATA and
RESTORE.
\\ \hline
REM string & REM allows remarks to be included in a program.
Note the Microsoft style use of the single quotation mark to denote
remarks is also supported and is preferred.
\\ \hline
RESTORE & Resets the line and position counters for DATA and READ statements to
the top of the program file.
\\ \hline
RETURN & RETURN concludes a subroutine called by GOSUB and returns to the
statement after the GOSUB.
\\ \hline
RUN [line] [file\$] & Executes the program in memory. If a line number is supplied then
execution will begin at that line, otherwise it will start at the beginning
of the program. Or, if a file name (file\$) is supplied, the current program
will be erased and that program will be loaded from the current drive and
executed. This enables one program to load and run another.
Example: RUN “TEST.BAS”
If an extension is not specified “.BAS” will be added to the file name.
\\ \hline
SEEK [\#]fnbr, pos & Will position the read/write pointer in a file that has been opened for
RANDOM access to the 'pos' byte.
The first byte in a file is numbered one so  SEEK \#5,1 will position the read/write pointer to the start of the file.
See Appendix I for more details on random file access.
\\ \hline
SUB xxx (arg1 [,arg2, …])
<statements>
<statements>
END SUB
& Defines a callable subroutine. This is the same as adding a new
command to MMBasic while it is running your program.
'xxx' is the subroutine name and it must meet the specifications for
naming a variable. 'arg1', 'arg2', etc are the arguments or parameters to
the subroutine.
Every definition must have one END SUB statement. When this is
reached the program will return to the next statement after the call to the
subroutine. The command EXIT SUB can be used for an early exit.
You use the subroutine by using its name and arguments in a program
just as you would a normal command. For example:  MySub a1, a2
When the subroutine is called each argument in the caller is matched to
the argument in the subroutine definition. These arguments are available
only inside the subroutine. Subroutines can be called with a variable
number of arguments. Any omitted arguments in the subroutine's list
will be set to zero or a null string.
Arguments in the caller's list that are a variable (ie, not an expression or
constant) will be passed by reference to the subroutine. This means that
any changes to the corresponding argument in the subroutine will also be
copied to the caller's variable and therefore may be accessed after the
subroutine has ended. Brackets around the argument list in both the caller and the definition are
optional.
\\ \hline
SYSTEM command\$ & Submit 'command\$' to the operating system.
It can be any command recognised by the command window in
Windows XP/Vista/7. The available commands are listed here:
http://ss64.com/nt
For example, this will set the window to blue lettering on a yellow
background: SYSTEM "COLOR 1E"
Note that the command is executed in a different instance of the
command processor to MMBasic so some commands (like "CD ..") will
have no effect.
\\ \hline
TROFF & Turns the trace facility off; see TRON.
\\ \hline
TRON & Turns on the trace facility. This facility will print the number of each line
(counting from the beginning of the program) in square brackets as the
program is executed. This is useful in debugging programs.
\\ \hline
ABS( number ) & Returns the absolute value of the argument 'number' (ie, any negative
sign is removed and the positive number is returned).
\\ \hline
ASC( string\$ ) & Returns the ASCII code for the first letter in the argument ‘string\$’.
\\ \hline
ATN( number ) & Returns the arctangent value of the argument 'number' in radians.
\\ \hline
BIN\$( number ) & Returns a string giving the binary (base 2) value for the 'number'.
\\ \hline
CHR\$( number ) & Returns a one-character string consisting of the character corresponding
to the ASCII code indicated by argument 'number'.
\\ \hline
CINT( number ) & Round numbers with fractional portions up or down to the next whole
number or integer.
For example,
45.47 will round to 45
45.57 will round to 46
-34.45 will round to -34
-34.55 will round to -35
See also INT() and FIX().
\\ \hline
CLR\$( )
or
CLR\$( fg )
or
CLR\$( fg, bg )
 & Returns a string containing embedded codes to select colours in a string.
'fg' is the foreground colour and 'bg' is the background colour. If no
parameters are specified both the foreground and background colours
will be reset to the defaults set by the last COLOUR command.
Example, this will display yellow letters on a red background:
PRINT CLR\$(YELLOW,RED) " ALARM "
This function simply generates a two character string where the first
character is the number 128 plus the foreground colour number and the
second character is the number 192 plus the background colour number.
\\ \hline
COS( number ) & Returns the cosine of the argument 'number' in radians.
\\ \hline
DEG( radians ) & Converts 'radians' to degrees.
\\ \hline
DIR\$( fspec, type )
or
DIR\$( fspec )
or
DIR\$( )
 & Will search the NOR flash for files and return the names of entries found.
'fspec' is a file specification using wildcards the same as used by the
FILES command. Eg, "*.*" will return all entries, "*.TXT" will return
text files.
'type' is the type of entry to return and can be one of:
VOL
Search for the volume label only
DIR
Search for directories only
FILE
Search for files only (the default if 'type' is not specified)
The function will return the first entry found. To retrieve subsequent
entries use the function with no arguments. ie, DIR\$( ). The return of
an empty string indicates that there are no more entries to retrieve.
This example will print all the files in a directory:
f\$ = DIR\$("*.*", FILE)
DO WHILE f\$ <> ""
PRINT f\$
f\$ = DIR\$()
LOOP
\\ \hline
EOF( [\#]nbr ) & Will return true if the file previously opened for INPUT with the file
number ‘nbr’ is positioned at the end of the file.
If used on a file number opened as a serial port this function will return
true if there are no characters waiting in the receive buffer.
The \# is optional. Also see the OPEN, INPUT and LINE INPUT
commands and the INPUT\$ function.
\\ \hline
EXP( number ) & Returns the exponential value of 'number'.
\\ \hline
FIX( number ) & Truncate a number to a whole number by eliminating the decimal point
and all characters to the right of the decimal point.
For example 9.89 will return 9 and -2.11 will return -2.

The major difference between FIX and INT is that FIX provides a true
integer function (ie, does not return the next lower number for negative
numbers as INT() does). This behaviour is for Microsoft compatibility.
See also CINT() .
\\ \hline
FORMAT\$( nbr [, fmt\$] ) & Will return a string representing ‘nbr’ formatted according to the
specifications in the string ‘fmt\$’.
The format specification starts with a % character and ends with a letter.
Anything outside of this construct is copied to the output as is.
The structure of a format specification is:
% [flags] [width] [.precision] type
Where ‘flags’ can be:
-
Left justify the value within a given field width
0
Use 0 for the pad character instead of space
+
Forces the + sign to be shown for positive numbers
space
Causes a positive value to display a space for the sign.
Negative values still show the – sign
‘width’ is the minimum number of characters to output, less than this the
number will be padded, more than this the width will be expanded.
‘precision’ specifies the number of fraction digits to generate with an e,
or f type or the maximum number of significant digits to generate with a
g type. If specified, the precision must be preceded by a dot (.).
‘type’ can be one of:
g
Automatically format the number for the best presentation.
f
Format the number with the decimal point and following
digits
e
Format the number in exponential format
If uppercase G or F is used the exponential output will use an uppercase
E. If the format specification is not specified “%g” is assumed.
Examples:
format\$(45) will return 45
format\$(45, “%g”) will return 45
format\$(24.1, “%g”) will return 24.1
format\$(24.1,”%f”) will return 24.100000
format\$(24.1, “%e”) will return 2.410000e+01
format\$(24.1,"%09.3f") will return 00024.100
format\$(24.1,"%+.3f") will return +24.100
format\$(24.1,"**%-9.3f**") will return **24.100  **
\\ \hline
HEX\$( number ) & Returns a string giving the hexadecimal (base 16) value for the 'number'.
\\ \hline
INKEY\$ & Checks the input of a character.
\\ \hline
INPUT\$(nbr, [\#]fnbr) & Will return a string composed of ‘nbr’ characters read from a file
previously opened for INPUT with the file number ‘fnbr’. This function
will read all characters including carriage return and new line without
translation.
When reading from a serial communications port this will return as
many characters as are waiting in the receive buffer up to ‘nbr’. If there
are no characters waiting it will immediately return with an empty string.
The \# is optional. Also see the OPEN command.
\\ \hline
INSTR( [start-position,]
string-searched\$, string-
pattern\$ )
 & Returns the position at which 'string-pattern\$' occurs in 'string-
searched\$', beginning at 'start-position'.
\\ \hline
INT( number ) & Truncate an expression to the next whole number less than or equal to
the argument. For example 9.89 will return 9 and -2.11 will return -3.
This behaviour is for Microsoft compatibility, the FIX() function
provides a true integer function.
See also CINT() .
\\ \hline
LEFT\$( string\$, number-
of-chars )
 & Returns a substring of ‘string\$’ with ‘number-of-chars’ from the left
(beginning) of the string.
\\ \hline
LEN( string\$ ) & Returns the number of characters in 'string\$'.
\\ \hline
LOC( [\#]fnbr ) & For a file opened as RANDOM this will return the current position of
the read/write pointer in the file. Note that the first byte in a file is
numbered 1. See Appendix I for more details of random file access.
For a serial port this will return the number of bytes received and
waiting in the receive buffer to be read.
The \# is optional.
\\ \hline
LOF( [\#]fnbr ) & For a file this will return the current length of the file in bytes.
For a serial port this will return the space (in characters) remaining in the
transmit buffer. Note that when the buffer is full MMBasic will pause
when adding a new character and wait for some space to become
available.
The \# is optional.
\\ \hline
LOG( number ) & Returns the natural logarithm of the argument 'number'. 
\\ \hline
LCASE\$( string\$ ) & Returns ‘string\$’ converted to lowercase characters.
\\ \hline
MID\$( string\$, start-
position-in-string[, number-
of-chars ] )
 & Returns a substring of ‘string\$’ beginning at ‘start-position-in-string’
and continuing for ‘number-of-chars’ bytes. If ‘number-of-chars’ is
omitted the returned string will extend to the end of ‘string\$’
\\ \hline
OCT\$( number ) & Returns a string giving the octal (base 8) representation of 'number'.
\\ \hline
PEEK( hiword, loword )
or
PEEK( keyword, +-offset)
or
PEEK(VAR var, +-offset)
& Will return a byte within the PIC32 virtual memory space.
The address is specifies by ‘hiword’ which is the top 16 bits of the
address while ‘loword’ is the bottom 16 bits.
Alternatively 'keyword' can be used and 'offset' is the +-offset from the
address of the keyword. The keyword can be 
PROGMEM (program
memory) or VARTBL (the variable table). 
You can also access the memory allocated to a variable by using the
variable's name ('var') preceded by the keyword VAR. This can be used
to access the individual bytes of a numeric variable or a large segment of
RAM allocated to an array (the first element of an array (eg, nbr(0)) is
the start of RAM allocated to the whole array).
See the POKE command for notes and warnings related to memory
access.
\\ \hline
PI & Returns the value of pi.
\\ \hline
POS & Returns the current cursor position in the line in characters.
\\ \hline
RAD( degrees ) & Converts 'degrees' to radians.
\\ \hline
RIGHT\$( string\$, number-
of-chars )
 & Returns a substring of ‘string\$’ with ‘number-of-chars’ from the right
(end) of the string.
\\ \hline
RND( number ) & Returns a pseudo-random number in the range of 0 to 0.999999. The
'number' value is ignored if supplied. The RANDOMIZE command
reseeds the random number generator.
\\ \hline
SGN( number ) & Returns the sign of the argument 'number', +1 for positive numbers, 0 for
0, and -1 for negative numbers.
\\ \hline
SIN( number ) & Returns the sine of the argument 'number' in radians.
\\ \hline
SPACE\$( number ) & Returns a string of blank spaces 'number' bytes long.
\\ \hline
SQR( number ) & Returns the square root of the argument 'number'.
\\ \hline
STR\$( number ) & Returns a string in the decimal (base 10) representation of 'number'.
\\ \hline
STRING\$( number, ascii-
value|string\$ )
 & Returns a string 'number' bytes long consisting of either the first
character of string\$ or the character representing the ASCII value ascii-
value.
\\ \hline
TAB( number ) & Outputs spaces until the column indicated by 'number' has been reached.
\\ \hline
TAN( number ) & Returns the tangent of the argument 'number' in radians.
\\ \hline
UCASE\$( string\$ ) & Returns ‘string\$’ converted to uppercase characters.
\\ \hline
VAL( string\$ ) & Returns the numerical value of the ‘string\$’. If 'string\$' is an invalid
number the function will return zero.
This function will recognise the \&H prefix for a hexadecimal number,
\&O for octal and \&B for binary.
\\ \hline
\end{tabular}
\end{table}

\begin{table}[]
\centering
\caption{Basic Instructions for System}
\label{Basic_Instructions_for_System}
\begin{tabular}{|p{4cm}|p{10cm}|}
\hline
\textbf{Basic Interpreter} & \textbf{Description}                                                             \\ \hline
3A & 3B 
\\ \hline
3A & 3B 
\\ \hline
3A & 3B 
\\ \hline
\end{tabular}
\end{table}


\begin{table}[]
\centering
\caption{Basic Instructions for File Operation}
\label{Basic_Instructions_for_File_Operation}
\begin{tabular}{|p{4cm}|p{10cm}|}
\hline
\textbf{Basic Interpreter} & \textbf{Description}                                                             \\ \hline
3A & 3B 
\\ \hline
3A & 3B 
\\ \hline
3A & 3B 
\\ \hline
\end{tabular}
\end{table}


\begin{table}[]
\centering
\caption{Basic Instructions for Communication}
\label{Basic_Instructions_for_Communication}
\begin{tabular}{|p{4cm}|p{10cm}|}
\hline
\textbf{Basic Interpreter} & \textbf{Description}                                                             \\ \hline
3A & 3B 
\\ \hline
3A & 3B 
\\ \hline
3A & 3B 
\\ \hline
\end{tabular}
\end{table}

\begin{table}[]
\centering
\caption{Basic Instructions for Numerical Operation}
\label{Basic_Instructions_for_Numerical_Operation}
\begin{tabular}{|p{4cm}|p{10cm}|}
\hline
\textbf{Basic Interpreter} & \textbf{Description}                                                             \\ \hline
3A & 3B 
\\ \hline
3A & 3B 
\\ \hline
3A & 3B 
\\ \hline
\end{tabular}
\end{table}

\begin{table}[]
\centering
\caption{Basic Instructions for String Operation}
\label{Basic_Instructions_for_String_Operation}
\begin{tabular}{|p{4cm}|p{10cm}|}
\hline
\textbf{Basic Interpreter} & \textbf{Description}                                                             \\ \hline
3A & 3B 
\\ \hline
3A & 3B 
\\ \hline
3A & 3B 
\\ \hline
\end{tabular}
\end{table}

\begin{table}[]
\centering
\caption{Basic Instructions for System}
\label{Basic_Instructions_for_System}
\begin{tabular}{|p{4cm}|p{10cm}|}
\hline
\textbf{Basic Interpreter} & \textbf{Description}                                                             \\ \hline
Memory & List the amount of memory currently in use. For example: \newline
 \textit{\newline 15kB (18\%) Program (528 lines) \newline 23kB (28\%) 52 Variables \newline 17kB (21\%) General \newline 28kB (33\%) Free} \newline \newline
Program memory is cleared by the NEW command. Variable and the general memory spaces are cleared by many commands (eg, NEW, RUN, LOAD, etc) as well as the specific commands CLEAR and ERASE. General memory is used by streams, file I/O buffers, etc.   
\\ \hline
\end{tabular}
\end{table}

%BASIC COMMANDS
%#	                Name	Tkn	Type
%   1)	              Memory	0080	C
%   2)	                   ?	0081	C
%   3)	               Clear	0082	C
%   4)	            Continue	0083	C
%   5)	                Data	0084	C
%   6)	                 Dim	0085	C
%   7)	                  Do	0086	C
%   8)	              ElseIf	0087	C
%   9)	             Else If	0088	C
%  10)	           Case Else	0089	C
%  11)	                Else	008a	C
%  12)	         Select Case	008b	C
%  13)	          End Select	008c	C
%  14)	                Case	008d	C
%  15)	               EndIf	008e	C
%  16)	              End If	008f	C
%  17)	        End Function	0090	C
%  18)	             End Sub	0091	C
%  19)	                 End	0092	C
%  20)	               Erase	0093	C
%  21)	               Error	0094	C
%  22)	            Exit For	0095	C
%  23)	            Exit Sub	0096	C
%  24)	       Exit Function	0097	C
%  25)	             Exit Do	0098	C
%  26)	                Exit	0099	C
%  27)	                 For	009a	C
%  28)	            Function	009b	C
%  29)	               GoSub	009c	C
%  30)	                GoTo	009d	C
%  31)	                  If	009e	C
%  32)	          Line Input	009f	C
%  33)	               Input	00a0	C
%  34)	                 Let	00a1	C
%  35)	                List	00a2	C
%  36)	               Local	00a3	C
%  37)	                Loop	00a4	C
%  38)	               Chain	00a5	C
%  39)	                Next	00a6	C
%  40)	                  On	00a7	C
%  41)	              Option	00a8	C
%  42)	               Print	00a9	C
%  43)	                Read	00aa	C
%  44)	                 Rem	00ab	C
%  45)	             Restore	00ac	C
%  46)	              Return	00ad	C
%  47)	                 Run	00ae	C
%  48)	                 Sub	00af	C
%  49)	               TROFF	00b0	C
%  50)	                TRON	00b1	C
%  51)	               Trace	00b2	C
%  52)	                Wend	00b3	C
%  53)	               While	00b4	C
%  54)	               Const	00b5	C
%  55)	                Open	00b6	C
%  56)	               Close	00b7	C
%  57)	               Files	00b8	C
%  58)	                Kill	00b9	C
%  59)	                Copy	00ba	C
%  60)	                Name	00bb	C
%  61)	                Seek	00bc	C
%  62)	             Library	00bd	C
%  63)	               Pause	00be	C
%  64)	               S0Out	00bf	C
%  65)	        SetVarAccess	00c0	C
%  66)	          SetTimeout	00c1	C
%  67)	             SYS.Set	00c2	C
%  68)	              SYSTEM	00c3	C
%  69)	            Dispatch	00c4	C
%  70)	                  @(	00c5	C
%  71)	              SYSLOG	00c6	C
%  72)	                 WDT	00c7	C
%  73)	              Memory	00c8	C
%  74)	           Copyright	00c9	C
%BASIC TOKENS
%#	                Name	Tkn	Type
%   1)	                Abs(	0080	fF
%   2)	               Acos(	0081	fF
%   3)	                Asc(	0082	iF
%   4)	               Asin(	0083	iF
%   5)	                Atn(	0084	fF
%   6)	               Bin$(	0085	sF
%   7)	               Chr$(	0086	sF
%   8)	               Cint(	0087	iF
%   9)	                Cos(	0088	fF
%  10)	                Deg(	0089	fF
%  11)	                Exp(	008a	fF
%  12)	                Fix(	008b	iF
%  13)	            Format$(	008c	sF
%  14)	               Hex$(	008d	sF
%  15)	              Inkey$	008e	sV
%  16)	              Instr(	008f	iF
%  17)	                Int(	0090	iF
%  18)	             LCase$(	0091	sF
%  19)	              Left$(	0092	sF
%  20)	             Ltrim$(	0093	sF
%  21)	                Len(	0094	iF
%  22)	                Log(	0095	fF
%  23)	               Mid$(	0096	sF
%  24)	         MM.CmdLine$	0097	sV
%  25)	              MM.Ver	0098	fV
%  26)	               Oct$(	0099	sF
%  27)	               Peek(	009a	iF
%  28)	                 Pi(	009b	fF
%  29)	                Pos(	009c	iF
%  30)	                Rad(	009d	fF
%  31)	             Right$(	009e	sF
%  32)	                Rnd(	009f	fF
%  33)	                 Rnd	00a0	fV
%  34)	              Round(	00a1	fF
%  35)	             Rtrim$(	00a2	sF
%  36)	                Sgn(	00a3	iF
%  37)	                Sin(	00a4	fF
%  38)	             Space$(	00a5	sF
%  39)	                Spc(	00a6	sF
%  40)	                Sqr(	00a7	fF
%  41)	               Str$(	00a8	sF
%  42)	            String$(	00a9	sF
%  43)	                Tab(	00aa	sF
%  44)	                Tan(	00ab	fF
%  45)	             UCase$(	00ac	sF
%  46)	                Val(	00ad	fF
%  47)	                 For	00ae	
%  48)	                Else	00af	
%  49)	               GoSub	00b0	
%  50)	                GoTo	00b1	
%  51)	                Step	00b2	
%  52)	                Then	00b3	
%  53)	                  To	00b4	
%  54)	               Until	00b5	
%  55)	               While	00b6	
%  56)	                   ^	00b7	fO
%  57)	                   *	00b8	fO
%  58)	                   /	00b9	fO
%  59)	                   \	00ba	iO
%  60)	                 Mod	00bb	iO
%  61)	                   +	00bc	fO
%  62)	                   -	00bd	fO
%  63)	                 Not	00be	fO
%  64)	                  <<	00bf	iO
%  65)	                  >>	00c0	iO
%  66)	                  <>	00c1	fO
%  67)	                  >=	00c2	fO
%  68)	                  =>	00c3	fO
%  69)	                  <=	00c4	fO
%  70)	                  =<	00c5	fO
%  71)	                   <	00c6	fO
%  72)	                   >	00c7	fO
%  73)	                   =	00c8	fO
%  74)	                 And	00c9	iO
%  75)	                  Or	00ca	iO
%  76)	                 Xor	00cb	iO
%  77)	                   ~	00cc	iO
%  78)	                Eof(	00cd	iF
%  79)	                Loc(	00ce	iF
%  80)	                Lof(	00cf	iF
%  81)	                  As	00d0	
%  82)	            MM.Errno	00d1	iV
%  83)	             Input$(	00d2	sF
%  84)	           MM.Fname$	00d3	sV
%  85)	               Dir$(	00d4	sF
%  86)	                Abs(	00d5	fF
%  87)	           Unixtime(	00d6	iF
%  88)	              Ticks(	00d7	iF
%  89)	         Timestamp$(	00d8	sF
%  90)	              Date$(	00d9	sF
%  91)	            Weekday(	00da	iF
%  92)	            Yearday(	00db	iF
%  93)	       CheckTimeout(	00dc	iF
%  94)	            SYS.Get(	00dd	fF
%  95)	               S0In(	00de	iF
%  96)	        Temperature(	00df	iF
%  97)	         RS485Read$(	00e0	sF
%  98)	       RS485ReadLn$(	00e1	sF
%  99)	          RS485Read(	00e2	iF
% 100)	         RS485Write(	00e3	iF
% 101)	       RS485ReadEx$(	00e4	sF
% 102)	     RS485ReadLnEx$(	00e5	sF
% 103)	        RS485ReadEx(	00e6	iF
% 104)	       RS485WriteEx(	00e7	iF
% 105)	            D0Start(	00e8	iF
% 106)	          D0ReadLn$(	00e9	sF
% 107)	              D0End(	00ea	iF
% 108)	            D0Read$(	00eb	sF
% 109)	             D0Read(	00ec	iF
% 110)	            D0Write(	00ed	iF
% 111)	       SocketServer(	00ee	iF
% 112)	       SocketClient(	00ef	iF
% 113)	    SocketConnected(	00f0	iF
% 114)	         SocketRead(	00f1	iF
% 115)	        SocketRead$(	00f2	sF
% 116)	      SocketReadLn$(	00f3	sF
% 117)	        SocketWrite(	00f4	iF
% 118)	        SocketClose(	00f5	iF
% 119)	       SocketOption(	00f6	iF
% 120)	               CRC$(	00f7	sF
% 121)	             Split$(	00f8	sF
% 122)	      StreamSearch$(	00f9	sF
% 123)	     EnoceanReceive(	00fa	iF
% 124)	    EnoceanTransmit(	00fb	iF
% 125)	        MQTTConnect(	00fc	iF
% 126)	     MQTTConnectTLS(	00fd	iF
% 127)	     MQTTDisconnect(	00fe	iF
% 128)	        MQTTPublish(	00ff	iF
% 129)	      MQTTSubscribe(	0180	iF
% 130)	    MQTTUnsubscribe(	0181	iF
% 131)	   MQTTSubscription(	0182	iF
% 132)	               Conv(	0183	fF
% 133)	                  @(	0184	fF
% 134)	           SetTimer(	0185	iF
% 135)	          KillTimer(	0186	iF
% 136)	         CrontabAdd(	0187	iF
% 137)	      CrontabRemove(	0188	iF
% 138)	                  As	0189	

\chapter{Specification}
% TODO measure weight and dimensions
Matrix with features
horizontal	EMDO101	EMDO102	EMDO103 MBUS	EMDO104 KNX
USB Power Supply o - - -
12V/24V Power Suply - o o o
D0 Interface o - - -
RS485 1 2 - -
MBUS - - o -
KNX - - - o
S0 Electrical Outputs 2 2 2 2 
S0 Electrical Inputs 2 2 2 2
S0 Optical Input 1 - - - -
Relay Outputs 1 2 2 2
1 Wire - 1 1 1 
Ethernet 1 1 1 1
Power Consumption < 1Watt (all)
Dimensions TBD
Weight TBD

\chapter{Credits and Licenses}
pls extract from credit.txt





%----------------------------------------------------------------------------------------
%	PART
%----------------------------------------------------------------------------------------

\part{Parte Dos}

%----------------------------------------------------------------------------------------
%	CHAPTER 3
%----------------------------------------------------------------------------------------



%----------------

%----------------------------------------------------------------------------------------
%	BIBLIOGRAPHY
%----------------------------------------------------------------------------------------

\chapter*{Bibliografía}
\addcontentsline{toc}{chapter}{\textcolor{ocre}{Bibliografía}}
\section*{Books}
\addcontentsline{toc}{section}{Books}
\printbibliography[heading=bibempty,type=book]

\begin{itemize}
	\item GREENE, W.H. (2003) “Econometric Analysis”5ª edición. Prentice Hall N.J. Capítulo 21
\\\\
    \item WOOLDRIDGE, J.M. (2010) “Introducción a la Econometría: Un Enfoque Moderno". 4ª edición. Cengage Learning. Capítulo 17

\end{itemize}


%----------------------------------------------------------------------------------------
%	INDEX
%----------------------------------------------------------------------------------------

\cleardoublepage
\phantomsection
\setlength{\columnsep}{0.75cm}
\addcontentsline{toc}{chapter}{\textcolor{ocre}{Índice Alfabético}}
\printindex

%----------------------------------------------------------------------------------------

\end{document}